\paragrafos{Άξονας συμμετρίας σχήματος}{Β.2.2, «Άξονας συμμετρίας»}

\begin{anagke}
Η Α3 διαλέγει τα σχήματα για τα παρτέρια του κήπου: τετράγωνο, ορθογώνιο, κύκλο, ισόπλευρο τρίγωνο και ένα παραλληλόγραμμο. Ο Βασίλης προτείνει να τα κόψουν από χαρτόνι και να τα διπλώσουν. «Αν διπλώσω ένα τετράγωνο στη μέση, συμπίπτουν τα δύο μισά», λέει. «Αλλά μπορώ να το διπλώσω και με άλλον τρόπο: από γωνία σε γωνία, και πάλι συμπίπτουν!». Η Αγγελική μετράει: το τετράγωνο διπλώνεται με τέσσερις διαφορετικούς τρόπους. Το ορθογώνιο με δύο. Το παραλληλόγραμμο καθόλου: όσο και αν προσπαθήσουν, τα δύο μισά δεν συμπίπτουν.

«Και ο κύκλος;» ρωτάει ο Νίκος. «Μπορώ να τον διπλώσω όπως θέλω, αρκεί η δίπλωση να περνά από το κέντρο!». Η Πολυξένη το γράφει: «Κάθε διπλωμένη ευθεία είναι ένας άξονας συμμετρίας του σχήματος». Η Ρήνη βρίσκει ότι ένα ισοσκελές τρίγωνο έχει έναν άξονα και ένα ισόπλευρο τρεις.

Στην παράγραφο αυτή ορίζουμε τον άξονα συμμετρίας ενός σχήματος και μετράμε πόσους άξονες έχουν τα γνωστά μας σχήματα.
\end{anagke}

\begin{lexi}{άξονας συμμετρίας σχήματος}{άξονας συμμετρίας σχήματος}
Ένα σχήμα έχει \lexiword{άξονα συμμετρίας} όταν υπάρχει ευθεία που το χωρίζει σε δύο μέρη που εφαρμόζουν τέλεια όταν διπλώσουμε κατά μήκος της: το σχήμα είναι δηλαδή το δικό του καθρέφτισμα.
\end{lexi}

\begin{metafrasi}
Η ευθεία που διπλώνει το σχήμα σε δύο ίσα μισά. & άξονας συμμετρίας\\
Το συμμετρικό του σχήματος ως προς τον άξονά του. & το ίδιο το σχήμα\\
Τετράγωνο. & $4$ άξονες\\
Ορθογώνιο (όχι τετράγωνο). & $2$ άξονες\\
Κύκλος. & άπειροι άξονες (κάθε διάμετρος)\\
\end{metafrasi}

\tmima{Η θεωρία}

\begin{orismos}{άξονας συμμετρίας}
\textbf{Άξονας συμμετρίας} ενός σχήματος είναι η ευθεία που το χωρίζει σε δύο μέρη τα οποία συμπίπτουν όταν διπλώσουμε το σχήμα κατά μήκος της. Τότε το συμμετρικό του σχήματος ως προς αυτήν είναι το ίδιο το σχήμα.
\end{orismos}

\begin{center}
\begin{tikzpicture}[scale=0.8]
\begin{scope}
\draw[thick] (0,0) rectangle (2,2); \draw[gblue,dashed] (1,-0.3)--(1,2.3) (-0.3,1)--(2.3,1) (-0.3,-0.3)--(2.3,2.3) (-0.3,2.3)--(2.3,-0.3);
\node at (1,-0.8) {\small τετράγωνο: $4$};
\end{scope}
\begin{scope}[xshift=4cm]
\draw[thick] (0,0) rectangle (3,1.8); \draw[gblue,dashed] (1.5,-0.3)--(1.5,2.1) (-0.3,0.9)--(3.3,0.9);
\node at (1.5,-0.8) {\small ορθογώνιο: $2$};
\end{scope}
\begin{scope}[xshift=9cm]
\draw[thick] (0,0)--(2.4,0)--(1.2,2.08)--cycle; \draw[gblue,dashed] (1.2,-0.3)--(1.2,2.4) (0,0)--(1.8,1.56) (2.4,0)--(0.6,1.56);
\node at (1.2,-0.8) {\small ισόπλευρο: $3$};
\end{scope}
\end{tikzpicture}
\end{center}

Μερικά παραδείγματα:
\begin{itemize}[leftmargin=1.2em,itemsep=0.3ex,topsep=2pt]
\item \textbf{Κύκλος}: κάθε ευθεία που περνά από το κέντρο είναι άξονας συμμετρίας (άπειροι).
\item \textbf{Τετράγωνο} $4$ (οι δύο μεσοκάθετοι των πλευρών και οι δύο διαγώνιοι). \textbf{Ορθογώνιο} $2$ (οι μεσοκάθετοι των πλευρών). \textbf{Ρόμβος} $2$ (οι διαγώνιοι).
\item \textbf{Παραλληλόγραμμο} (που δεν είναι ορθογώνιο ή ρόμβος): κανένας.
\item \textbf{Ισοσκελές τρίγωνο}: $1$, η ευθεία από την κορυφή που διχοτομεί τη γωνία της. \textbf{Ισόπλευρο}: $3$. \textbf{Σκαληνό}: κανένας.
\item \textbf{Κανονικό $\nu$-γωνο}: $\nu$ άξονες.
\item \textbf{Ευθύγραμμο τμήμα}: $2$ (η ευθεία του και η μεσοκάθετός του). \textbf{Γωνία}: ο φορέας της διχοτόμου της.
\end{itemize}

\begin{prosoxi}
\begin{itemize}[leftmargin=1.2em,itemsep=0.3ex,topsep=2pt]
\item Για να είναι άξονας συμμετρίας δεν αρκεί να χωρίζει το σχήμα σε δύο ίσα εμβαδά: πρέπει τα δύο μέρη να συμπίπτουν. Η διαγώνιος ενός (μη τετράγωνου) ορθογωνίου το χωρίζει σε δύο ίσα εμβαδά αλλά δεν είναι άξονας.
\item Το παραλληλόγραμμο δεν έχει άξονα συμμετρίας (παρότι έχει κέντρο συμμετρίας, όπως θα δούμε).
\item Ο άξονας συμμετρίας είναι ολόκληρη \emph{ευθεία}, όχι μόνο το κομμάτι μέσα στο σχήμα.
\item Ένα σχήμα μπορεί να έχει πολλούς άξονες συμμετρίας, και όλοι διέρχονται από το ίδιο σημείο όταν είναι πάνω από έναν.
\end{itemize}
\end{prosoxi}

\tmima{Λυμένα παραδείγματα}

\paradeigma{Τετράγωνο και ορθογώνιο}
\tixero Ένα τετράγωνο και ένα ορθογώνιο $4\times6$.\par
\tithelo Πόσους άξονες έχει το καθένα.\par
\possindeo Το τετράγωνο διπλώνεται: στη μέση οριζόντια, στη μέση κατακόρυφα και κατά μήκος κάθε διαγωνίου: $4$. Το ορθογώνιο μόνο οριζόντια και κατακόρυφα στη μέση: $2$ (οι διαγώνιοι δεν δουλεύουν).\par
\apantisi $4$ και $2$.\par
\elegxos Αν διπλώσω το ορθογώνιο στη διαγώνιο, μια γωνία «βγαίνει έξω».

\paradeigma{Ο κύκλος}
\tixero Ο κύκλος $(\mathrm{Ο},\rho)$.\par
\tithelo Οι άξονες συμμετρίας.\par
\possindeo Κάθε διάμετρος $\delta$ χωρίζει τον κύκλο σε δύο ημικύκλια που εφαρμόζουν. Η ευθεία της διαμέτρου είναι άξονας.\par
\apantisi Άπειροι: οι ευθείες που περνούν από το $\mathrm{Ο}$.\par
\elegxos Δεν υπάρχει άξονας που να μην περνά από το κέντρο.

\paradeigma{Τρίγωνα}
\tixero Ισόπλευρο, ισοσκελές, σκαληνό τρίγωνο.\par
\tithelo Οι άξονες συμμετρίας.\par
\possindeo Το ισόπλευρο έχει από κάθε κορυφή έναν άξονα (στη μέση της απέναντι πλευράς): $3$. Το ισοσκελές έχει μόνο τον άξονα της κορυφής: $1$. Το σκαληνό: κανένα.\par
\apantisi $3$, $1$, $0$.\par
\elegxos Στο ισοσκελές, οι δύο ίσες πλευρές είναι συμμετρικές.

\paradeigma{Τα γράμματα}
\tixero Τα γράμματα $\mathrm{Α}$, $\mathrm{Β}$, $\mathrm{Η}$, $\mathrm{Ο}$, $\mathrm{Ζ}$.\par
\tithelo Ποια έχουν άξονα συμμετρίας.\par
\possindeo Το $\mathrm{Α}$ έχει κατακόρυφο άξονα. Το $\mathrm{Β}$ έχει οριζόντιο. Το $\mathrm{Η}$ έχει δύο (οριζόντιο και κατακόρυφο). Το $\mathrm{Ο}$ έχει δύο (αν είναι έλλειψη) ή άπειρους (αν κύκλος). Το $\mathrm{Ζ}$ κανέναν.\par
\apantisi Όλα εκτός από το $\mathrm{Ζ}$.\par
\elegxos Το $\mathrm{Ζ}$ έχει όμως κέντρο συμμετρίας.

\tmima{Ασκήσεις}

\zoni{Το κατάλαβα}{πόσοι άξονες}

\ex Κόψε από χαρτί ένα τετράγωνο και βρες όλους τους άξονες συμμετρίας του με δίπλωση.

\ex Πόσους άξονες συμμετρίας έχει: α) το ορθογώνιο ($4\times6$) \quad β) ο ρόμβος \quad γ) το ισόπλευρο τρίγωνο \quad δ) το ισοσκελές \quad ε) το σκαληνό \quad στ) το παραλληλόγραμμο.

\ex Πόσους άξονες συμμετρίας έχει το κανονικό πεντάγωνο, το κανονικό εξάγωνο και ο κύκλος;

\ex Σχεδίασε ένα ευθύγραμμο τμήμα $\mathrm{ΑΒ}$. Βρες τους άξονες συμμετρίας του.

\ex Σχεδίασε μια γωνία $x\mathrm{Ο}y$. Βρες τον άξονα συμμετρίας της.

\ex Ποια από τα κεφαλαία γράμματα του ελληνικού αλφαβήτου $\mathrm{Α}$, $\mathrm{Β}$, $\Gamma$, $\Delta$, $\mathrm{Ε}$, $\mathrm{Η}$, $\mathrm{Θ}$, $\mathrm{Ι}$, $\mathrm{Κ}$, $\Lambda$, $\mathrm{Μ}$, $\mathrm{Ν}$, $\mathrm{Ο}$, $\mathrm{Π}$, $\mathrm{Τ}$, $\mathrm{Φ}$, $\mathrm{Χ}$, $\Psi$ έχουν άξονα συμμετρίας;

\ex Σωστό ή λάθος; α) Κάθε σχήμα έχει άξονα συμμετρίας. \quad β) Ο κύκλος έχει άπειρους άξονες. \quad γ) Η διαγώνιος ορθογωνίου είναι άξονας συμμετρίας του. \quad δ) Το τετράγωνο έχει $4$ άξονες. \quad ε) Το ισοσκελές τρίγωνο έχει τουλάχιστον έναν άξονα.

\zoni{Το χρησιμοποιώ}{κήπος, σχέδια, σημαίες}

\ex \textbf{Το παρτέρι.} Ένα παρτέρι είναι κανονικό εξάγωνο. Πόσοι είναι οι άξονες συμμετρίας του; Πόσοι περνούν από δύο κορυφές;

\ex \textbf{Η πόρτα.} Η πόρτα του σχολείου έχει σχήμα ορθογωνίου διαστάσεων $1{,}2\times2{,}1$ μ. Πόσους άξονες συμμετρίας έχει το σχήμα της;

\ex \textbf{Το συντριβάνι.} Σε ένα ορθογώνιο παρτέρι $4\times6$ μ. θέλουμε να βάλουμε ένα συντριβάνι. Σε ποια σημεία θα είναι «συμμετρικό» ως προς όλους τους άξονες;

\ex \textbf{Το λογότυπο.} Σχεδίασε ένα λογότυπο για την Α3 που να έχει ακριβώς δύο άξονες συμμετρίας.

\ex \textbf{Ο χαρταετός.} Ένας χαρταετός έχει σχήμα ρόμβου. Πόσους άξονες συμμετρίας έχει; Πού τέμνονται;

\zoni{Το σκέφτομαι}{όταν η δίπλωση δεν δουλεύει}

\ex Γιατί η διαγώνιος ενός ορθογωνίου (που δεν είναι τετράγωνο) δεν είναι άξονας συμμετρίας, παρότι το χωρίζει σε δύο ίσα τρίγωνα;

\ex Μπορεί ένα σχήμα να έχει ακριβώς δύο άξονες συμμετρίας; Ακριβώς τρεις; Δώσε παραδείγματα.

\ex Αν ένα σχήμα έχει δύο άξονες που τέμνονται, πού βρίσκεται ένα σημείο που «ανήκει» και στους δύο;

\ex Γιατί ο άξονας συμμετρίας ενός ισοσκελούς τριγώνου διχοτομεί τη γωνία της κορυφής;

\zoni{Για το σπίτι}{συμμετρίες γύρω σου}

\ex Βρες πέντε σχήματα στο σπίτι με έναν άξονα, τρία με δύο άξονες και ένα με τρεις ή περισσότερους.

\ex Δίπλωσε ένα χαρτί και κόψε ένα σχήμα. Άνοιξέ το και βρες τον άξονα συμμετρίας του.
